\section{Discussion}
\label{sec:disc}

\subsection{What to claim}
The defensible claims are architectural and methodological: (1)~ECNetBench provides a leakage-aware multi-task enterprise benchmark with multi-seed frozen instances; (2)~ECN-v3 closes the loop from detection to healing \emph{decision support} with TreeSHAP RCA; (3)~leakage-safe temporal enrichment is the primary driver of the T1 lift from v2 to the final head; (4)~anchored fusion is preferred over stacking on this benchmark (higher mean AUPRC, lower complexity); (5)~T2 remains best served by a telem logistic head under extreme rarity.

\subsection{What not to claim}
We do not claim statistically significant superiority over random forest on T1 ($p{=}0.688$), significant superiority of anchored over stacking ($p{=}0.375$), live deployment on physical AOS-CX switches, or that multi-agent fusion alone beats strong tabular detectors without feature enrichment. TabNet and true GraphSAGE likewise do not overturn the evidence-selected heads on this protocol. Impact AUPRC of $1.0$ for some models on this synthetic label construction is reported for completeness but interpreted cautiously as potential label/feature collinearity rather than solved service-impact prediction in the wild.

\subsection{Practical NetOps implications}
Operators can use ECN as a sandbox for ranking alerts, proposing RCA categories with SHAP attributions, and suggesting remediation classes before human approval. Prefer the anchored head for T1 triage; prefer telem logistic scores for T2 horizon ranking; apply optional Beta calibration when displaying probabilities.

\subsection{Governance and operational implications}
Because ECN-v3 closes a multi-agent loop that can recommend remediation classes, deployment should retain human approval gates for delegated actions, consistent with governance analyses of agentic systems and agent-to-agent delegation~\cite{rah2026agenticreview,rah2026masreview,rah2026a2a}. Leakage-safe temporal evaluation and checksummed instances support operational-data auditability~\cite{rah2026auditing}, while cybersecurity risk-control and zero-exposure guidance caution against exposing privileged recovery pathways through automated tooling~\cite{rah2026cybercontrols,rah2026zeroexposure}. If generative or retrieval-augmented assistants are later attached to NetOps workflows, enterprise generative-AI security and agentic-AI governance frameworks provide complementary controls beyond detector metrics~\cite{rah2026entgenai,rah2026raggovernance}. Semantic tracing practices for LLM-based multi-agent stacks offer a useful template for governing agent decisions alongside TreeSHAP RCA~\cite{rah2026semantictracing}. Digital-transformation studies further remind operators that legacy constraints and human resistance often dominate outcomes even when models improve ranking scores~\cite{rah2026dtbeyond,rah2026dtfails,rah2026govresistance,rah2026govdtcyber}.
